\chapter{Research Objectives}
\label{sec:research_objectives}

The primary aim of this PhD project is to enhance informed decision-making in the diagnosis of ovarian cancer by improving the \gls{ADNEX} model, ensuring that it can be used safely and effectively across diverse clinical settings. The project integrates both applied (clinical) and methodological aspects, which must be clearly distinguished. Although the main focus is on the \gls{ADNEX} model in ovarian cancer, the methodological work is conducted in such a way that the insights of this PhD can be applied generally to the clinical prediction model field.

On the applied side, the aim is to evaluate, update, and externally validate the \gls{ADNEX} model to minimise performance heterogeneity and improve generalisability. The first step is to systematically review and meta-analyse \gls{ADNEX} model performance across settings, identifying sources of heterogeneity such as hospital type and examiner expertise (Chapters~\ref{sec:chp_ADNEXSR} and~\ref{sec:chp_ADNEXvsRMI}). Building on these findings, the next steps are to update the \gls{ADNEX} model with the insights obtained during the ten years that the model has been used in clinic and with more recent data that also includes patients that are not operated, to externally validate the updated model using new \gls{IOTA} studies employing techniques that quantify performance heterogeneity, and, if required, to develop a new prediction model with an optimised set of predictors.

On the methodological side, the aim is to advance approaches for evaluating and improving clinical prediction models, with a primary focus on accurate predicted probabilities, also known as model calibration, in clustered settings. A first objective is to investigate the impact of different algorithms and hyperparameter settings on predicted probabilities (Chapters~\ref{sec:chp_RFOverfitting} and~\ref{sec:chp_FundamentalProblem}). A second objective is to develop methods for calibration plots that integrate clustering information, enabling improved calibration assessment in studies with clustered data (Chapter~\ref{sec:chp_ClusteredCalibration}). A third objective is to explore approaches for incorporating clinical utility into model development, including designing a variable selection algorithm that prioritises clinical usefulness over purely statistical fit.

The expected outcome is a robust, well-calibrated, and clinically useful prediction model for ovarian cancer diagnosis with demonstrated generalisability, supporting informed decision-making in ovarian cancer management across different centres and examiners.
